\chapter{\texorpdfstring{The Systemic Form: Derivazione di
\(T_s^{(2),\text{sys}}\)}{The Systemic Form: Derivazione di T\_s\^{}\{(2),\textbackslash text\{sys\}\}}}\label{la-forma-sistemica-followszione-di-t_s2textsys}

\begin{quote}
\textbf{Core argument:} la stabilità di un sistema non è la media
delle stabilità dei suoi attori --- è quella media \emph{modulata dalla
phase coherence}, che amplifica l'antifragility collettiva e divide
l'involution collettiva. La differenza tra \(T_s^{(2)}\) local e
\(T_s^{(2),\text{sys}}\) sistemico misura quanto un sistema
\emph{spreca} per decoherence interna.
\end{quote}



\section{Stato del Problema}\label{stato-del-problema}

Il Ch.~3 ha followsto \(T_s^{(2)}\) come quantità \emph{local} di un
attore con componenti perfettamente allineate (\(\Phi = 1\)). Il \S{}3.1bis
ha introdotto il order parameter di Kuramoto
\(\Phi(t) = \left|\frac{1}{N}\sum_j e^{i\theta_j(t)}\right|\) e
l'transposition operator
\(\langle X \rangle_\Phi := \Phi \cdot \frac{1}{N}\sum_j X_j\).

Resta aperta la domanda: \textbf{come \(A\) e \(\Omega\) si trasformano
nel passaggio dalla scala local alla scala sistemica?} La tesi di
questo chapter è che la trasformazione è \emph{asimmetrica}: \(\Phi\)
si comporta come un \textbf{moltiplicatore} per l'antifragility (la
coerenza estende il guadagno da stress all'intero sistema) e come un
\textbf{divisore} per l'involution (la decoherence moltiplica gli
sprechi). Da qui la forma:

\[
A_{\text{sys}} = \Phi \cdot \langle A \rangle, \qquad \Omega_{\text{sys}} = \frac{\langle \Omega \rangle}{\Phi}
\]

Il resto del chapter follows, motiva, verification e interpreta questa
asimmetria.



\section{Setup Formale: il Campo dei Payoff
Locali}\label{setup-formal-il-campo-dei-payoff-locali}

Il sistema sociale è una popolazione di \(N\) accoppiamenti elementari
(cittadino-discorso, attore-egemone, élite-massa, impresa-mercato).
Ciascuno ha:

\begin{itemize}
\tightlist
\item
  una \textbf{fase} \(\theta_j(t) \in [0, 2\pi)\) --- posizione nel
  ciclo di legittimazione/decoherence dell'coupling;
\item
  una \textbf{natural frequency} \(\omega_j\) --- ritmo intrinseco
  dell'coupling (cicli elettorali, di mercato, mediatici);
\item
  un \textbf{instantaneous payoff}
  \(\Pi_j(t) = M_{e,j} \cdot Q_{0,j} \cdot (1 - \delta_{R,j} z_{R,j})\)
  ereditato dalla legge micro della Topografia (\S{}3.1.2);
\item
  un'\textbf{antifragility local}
  \(A_j = \partial^2 \Pi_j / \partial \sigma^2\);
\item
  un \textbf{coefficiente di involution local}
  \(\Omega_j = \Delta E_{\text{diss},j}/\Delta E_{\text{utili},j} - 1\).
\end{itemize}

\Hypoth{} \textbf{Ipotesi di mean field.} Trattiamo le fasi come
accoppiate via la dinamica di Kuramoto (1984):

\[
\frac{d\theta_j}{dt} = \omega_j + \frac{K}{N}\sum_{k=1}^{N} \sin(\theta_k - \theta_j)
\]

where \(K \geq 0\) è la \emph{forza di coupling sistemica} (densità
mediatica, integrazione istituzionale, infrastruttura comunicativa).
L'equation si riduce, in termini del order parameter \(\Phi\), a:

\[
\frac{d\theta_j}{dt} = \omega_j + K\Phi \sin(\psi - \theta_j)
\]

Esiste una critical threshold \(K_c = 2/(\pi g(0))\), where \(g(\omega)\) è
la distribuzione delle frequenze naturali, sotto cui \(\Phi = 0\)
(decoherence completa) e sopra cui \(\Phi > 0\) (synchronization
parziale).

\Proved{} \textbf{Confollowsnza formal:} \(\Phi\) non è un parameter
libero --- è determinato dalla densità infrastructural del sistema.
Sistemi con bassa \(K\) (frammentazione) hanno \(\Phi \to 0\) anche se
ciascun attore è perfettamente coerente con se stesso.



\section{Trasposizione di A: l'Antifragility
Sistemica}\label{trasposizione-di-a-lantifragilituxe0-sistemica}

\subsection{Argomento Strutturale}\label{argomento-structural}

Un attore antifragile (\(A_j > 0\)) trae beneficio individuale dalla
volatilità. Ma \emph{quel beneficio si trasferisce al sistema} solo se
il sistema è abbastanza coerente da \textbf{redistribuire} il guadagno:
condividere l'innovazione, propagare la lezione, integrare l'opzionalità
acquisita.

In un sistema decoerente, ogni attore antifragile cattura la propria
rendita; il sistema nel suo complesso non cresce da quel guadagno. In un
sistema coerente, lo stesso \(\langle A \rangle\) produce un guadagno
aggregato moltiplicato dalla capacità di trasferimento \(\Phi\).

\subsection{Forma Funzionale}\label{forma-funzionale}

\Hypoth{} \textbf{Postulato di trasposizione (forma più semplice non
banale):}

\[
A_{\text{sys}}(\Phi) = \Phi \cdot \langle A \rangle, \qquad \langle A \rangle = \frac{1}{N}\sum_j A_j
\]

Questa forma è la \textbf{first espansione di Taylor} di
\(A_{\text{sys}}\) attorno al regime di piena coerenza, soggetta alle
due condizioni al contorno:

\begin{itemize}
\tightlist
\item
  \(A_{\text{sys}}(\Phi = 1) = \langle A \rangle\) (la coerenza piena
  recupera la media);
\item
  \(A_{\text{sys}}(\Phi = 0) = 0\) (la decoherence piena annulla il
  trasferimento).
\end{itemize}

\Proved{} \textbf{Verifica empirica preliminare.} Confronto USA /
Brazil: gli USA hanno \(\langle A \rangle\) alto (Silicon Valley,
mercato dei capitali dinamico, ecosistema startup), ma \(\Phi \to\)
basso (polarizzazione, frammentazione del discorso, bassa coerenza
istituzionale). Il Brazil ha \(\langle A \rangle\) medio (jeitinho,
plasticità improvvisativa) ma \(\Phi\) alto (sincretismo culturale,
coerenza simbolica). La predizione DEQ$^2$ è
\(A_{\text{sys}}^{\text{USA}} \approx A_{\text{sys}}^{\text{BRA}}\) ---
coerente con il declino relativo dell'efficienza sistemica USA degli
anni 2020.

\Conj{} \textbf{Limite della forma lineare.} È plausibile che per
\(\Phi \to 1\) esistano effetti di saturation (diminishing returns)
che richiederebbero una forma
\(A_{\text{sys}} = (1 - e^{-\alpha\Phi})\langle A \rangle\). La forma
lineare è la più semplice consistente; alternative non lineari sono
rinviate a verification empirica nel paper EN.



\section{Trasposizione di $\Omega$: l'Involuzione
Sistemica}\label{trasposizione-di-ux3c9-linvolution-sistemica}

\subsection{Argomento Strutturale}\label{argomento-structural-1}

L'asimmetria con l'antifragility è essenziale. Mentre \(A\) produce
\emph{valore} che richiede coerenza per propagarsi, \(\Omega\) produce
\emph{spreco} che la decoherence \textbf{amplifica strutturalmente}.

In un sistema coerente, gli sprechi si compensano: la dissipazione di un
attore può alimentare l'utilità di un altro (l'energia residua del
settore A diventa input del settore B). In un sistema decoerente, ogni
attore dissipa indipendingmente, e gli sprechi si \textbf{sommano senza
compensazione} --- anzi si moltiplicano per duplicazione, competizione
improduttiva, perdita di standardizzazione.

\subsection{Forma Funzionale}\label{forma-funzionale-1}

\Hypoth{} \textbf{Postulato di trasposizione:}

\[
\Omega_{\text{sys}}(\Phi) = \frac{\langle \Omega \rangle}{\Phi}, \qquad \langle \Omega \rangle = \frac{1}{N}\sum_j \Omega_j
\]

Condizioni al contorno verificationte:

\begin{itemize}
\tightlist
\item
  \(\Omega_{\text{sys}}(\Phi = 1) = \langle \Omega \rangle\) (la
  coerenza piena recupera la media);
\item
  \(\Omega_{\text{sys}}(\Phi \to 0) \to \infty\) (la decoherence piena fa
  esplodere la dissipazione collettiva).
\end{itemize}

\Proved{} \textbf{Verifica empirica preliminare.} La China ha
\(\langle \Omega \rangle\) molto alto (Neijuan endemico nei settori 996,
sovracapacità industriale, \emph{involution} documentata da Xiang Biao)
ma \(\Phi\) alto (coerenza di comando del PCC).
\(\Omega_{\text{sys}}^{\text{CINA}} = \langle\Omega\rangle/\Phi \approx \langle\Omega\rangle\)
--- alto ma contenuto. L'India ha \(\langle \Omega \rangle\) medio ma
\(\Phi\) basso (frammentazione comunitaria, federalismo competitivo); la
sua \(\Omega_{\text{sys}}\) effettiva è proporzionalmente peggiore di
quanto suggerisca la metrica local. Coerente con la difficoltà indiana
a tradurre la propria base produttiva in coerenza geopolitica.

\Conj{} \textbf{Asimmetria notezionale: because $\Phi$ moltiplica A ma divide
$\Omega$?} È l'\textbf{asimmetria fondamentale tra valore e spreco}. Il valore
richiede \emph{trasporto} per moltiplicarsi (e \(\Phi\) è il
coefficiente di trasporto). Lo spreco non richiede trasporto --- si
moltiplica per duplicazione spontanea ogni volta che il trasporto è
assente. Questa asimmetria è la firma matematica della \textbf{seconda
legge della termodinamica applicata ai sistemi sociali}: l'ordine si
trasmette, il disordine si crea ovunque non sia inibito.



\section{\texorpdfstring{Derivazione di
\(T_s^{(2),\text{sys}}\)}{Derivazione di T\_s\^{}\{(2),\textbackslash text\{sys\}\}}}\label{followszione-di-t_s2textsys}

\subsection{Costruzione}\label{costruzione}

Partiamo dalla forma local del Ch.~3:

\[
T_s^{(2)} = \frac{s_{iii} \cdot z \cdot \delta \cdot (1 + A)}{\sigma \cdot \varepsilon_{\text{dis}} \cdot (1 + \Omega)}
\]

Per ottenere la forma sistemica, sostituiamo \(A \to A_{\text{sys}}\) e
\(\Omega \to \Omega_{\text{sys}}\) secondo i postulati di trasposizione,
mantenendo le altre variables come quantità sistemiche aggregate
(legitimacy collettiva \(s_{iii}\), dialectical complexity media \(z\),
orizzonte \(\delta\), volatilità ambientale \(\sigma\), coefficiente di
disimpegno \(\varepsilon_{\text{dis}}\)):

\[
\boxed{T_s^{(2),\text{sys}} = \frac{s_{iii} \cdot z \cdot \delta \cdot \left(1 + \Phi \langle A \rangle\right)}{\sigma \cdot \varepsilon_{\text{dis}} \cdot \left(1 + \dfrac{\langle \Omega \rangle}{\Phi}\right)}}
\]

\Proved{} \textbf{Coincide con la forma anticipata nell'Introduzione del
libro} (section \emph{Apparato formal chiave}). La followszione qui
presentata è la giustificazione formal di quella formula.

\subsection{Verifiche di Consistenza nei
Limiti}\label{verifiche-di-consistenza-nei-limiti}

{\def\LTcaptype{none} % do not increment counter
{\small\begin{longtable}[]{@{}
  >{\raggedright\arraybackslash}p{(\linewidth - 4\tabcolsep) * \real{0.3333}}
  >{\raggedright\arraybackslash}p{(\linewidth - 4\tabcolsep) * \real{0.3333}}
  >{\raggedright\arraybackslash}p{(\linewidth - 4\tabcolsep) * \real{0.3333}}@{}}
\toprule\noalign{}
\begin{minipage}[b]{\linewidth}\raggedright
Limite
\end{minipage} & \begin{minipage}[b]{\linewidth}\raggedright
Forma di \(T_s^{(2),\text{sys}}\)
\end{minipage} & \begin{minipage}[b]{\linewidth}\raggedright
Interpretazione
\end{minipage} \\
\midrule\noalign{}
\endhead
\bottomrule\noalign{}
\endlastfoot
\(\Phi \to 1\) &
\(\dfrac{s_{iii} z \delta (1+\langle A\rangle)}{\sigma \varepsilon_{\text{dis}} (1+\langle\Omega\rangle)}\)
& Recupero di \(T_s^{(2)}\) con \(A = \langle A\rangle\),
\(\Omega = \langle\Omega\rangle\) \Proved{} \\
\(\Phi \to 0\) &
\(\dfrac{s_{iii} z \delta}{\sigma \varepsilon_{\text{dis}}} \cdot \dfrac{1}{\infty} \to 0\)
& Decoerenza piena = collapse sistemico \Proved{} \\
\(\langle A\rangle \to \infty\), \(\Phi\) finito &
\(T_s^{(2),\text{sys}} \to \infty\) & Antifragility infinita salva il
sistema \Proved{} \\
\(\langle\Omega\rangle \to \infty\), \(\Phi\) finito &
\(T_s^{(2),\text{sys}} \to 0\) & Involuzione totale fa collassare il
sistema \Proved{} \\
\(K < K_c\) & \(\Phi = 0 \Rightarrow T_s^{(2),\text{sys}} = 0\) & Sotto
la soglia di coupling di Kuramoto, qualsiasi sistema collassa
\Proved{} \\
\end{longtable}}
}

\Proved{} \textbf{Risultato structural.} Tutti i limiti hanno
interpretazione coerente. Nessun limite produce comportamento
patologico.



\section{Il Punto Critico di $\Phi$}\label{il-punto-critico-di-ux3c6}

\subsection{La $\Phi$ Critica}\label{la-ux3c6-critica}

\Hypoth{} \textbf{Definizione operativa.} Dato un sistema con
\((s_{iii}, z, \delta, \sigma, \varepsilon_{\text{dis}}, \langle A\rangle, \langle\Omega\rangle)\)
fissati, la \textbf{critical threshold di coerenza} \(\Phi^*\) è il valore
for which \(T_s^{(2),\text{sys}}(\Phi^*) = 1\). Imponendo l'uguaglianza in
\(T_s^{(2),\text{sys}}\) e moltiplicando per \(\Phi^*\) entrambi i lati
si ottiene la forma \emph{quadratica esatta}:

\[
s_{iii} z \delta \langle A\rangle \cdot (\Phi^*)^2 + (s_{iii} z \delta - \sigma \varepsilon_{\text{dis}}) \cdot \Phi^* - \sigma \varepsilon_{\text{dis}} \langle\Omega\rangle = 0
\]

\Hypoth{} \textbf{Approssimazione lineare di servizio (per
\(\Phi^* \approx 1\)).} Trascurando il termine quadratico in \(\Phi^*\)
--- accettabile in regime di alta coerenza, vicino a 1 --- si ottiene la
forma \emph{linearizzata}:

\[
\Phi^*_{\text{lin}} \approx \frac{\sigma \varepsilon_{\text{dis}} \langle\Omega\rangle - s_{iii} z \delta}{s_{iii} z \delta \langle A\rangle - \sigma \varepsilon_{\text{dis}}}
\]

\Proved{} \textbf{Forma esatta.} La risoluzione completa della
quadratica via formula classica è followsta esplicitamente
nell'\textbf{Appendix A \S{}A.7} ed è da preferire per sistemi in regime
di forte decoherence (\(\Phi\) lontano da 1) --- caso paradigmatico, USA
del Ch.~7.

\Hypoth{} \textbf{Interpretazione.} Sistemi con \(\Phi(t) < \Phi^*\)
sono in \textbf{regime di phase transition imminente} (nel senso del
Ch.~3, \S{}3.5.3). Il monitoraggio empirico di \(\Phi\) diventa il primo
indicatore precoce di crisi sistemica.

\subsection{Le Quattro Posizioni
Strutturali}\label{le-quattro-posizioni-strutturali}

In funzione di \(\Phi\) e dei parameters locali, distinguiamo quattro
posizioni strutturali --- che mappano direttamente i quattro attori del
libro:

{\def\LTcaptype{none} % do not increment counter
{\small\begin{longtable}[]{@{}
  >{\raggedright\arraybackslash}p{(\linewidth - 8\tabcolsep) * \real{0.2000}}
  >{\raggedright\arraybackslash}p{(\linewidth - 8\tabcolsep) * \real{0.2000}}
  >{\raggedright\arraybackslash}p{(\linewidth - 8\tabcolsep) * \real{0.2000}}
  >{\raggedright\arraybackslash}p{(\linewidth - 8\tabcolsep) * \real{0.2000}}
  >{\raggedright\arraybackslash}p{(\linewidth - 8\tabcolsep) * \real{0.2000}}@{}}
\toprule\noalign{}
\begin{minipage}[b]{\linewidth}\raggedright
Posizione
\end{minipage} & \begin{minipage}[b]{\linewidth}\raggedright
\(\Phi\)
\end{minipage} & \begin{minipage}[b]{\linewidth}\raggedright
\(\langle A\rangle\)
\end{minipage} & \begin{minipage}[b]{\linewidth}\raggedright
\(\langle\Omega\rangle\)
\end{minipage} & \begin{minipage}[b]{\linewidth}\raggedright
Attore paradigmatico
\end{minipage} \\
\midrule\noalign{}
\endhead
\bottomrule\noalign{}
\endlastfoot
\textbf{Soliton metastabile} & n.a. (\(N=1\)) & molto alto & basso &
Musk \\
\textbf{Egemone decoerente} & basso & alto & medio & USA \\
\textbf{Rigido sincronizzato} & alto & basso & alto & China \\
\textbf{Bilanciere quantistico} & alto & medio & basso & Brazil \\
\end{longtable}}
}

\Proved{} \textbf{Risultato non banale.} Il \textbf{Bilanciere
quantistico} non è il sistema con i parameters locali migliori --- è
quello in cui \emph{la combinazione} di \(\Phi\) alto,
\(\langle A\rangle\) medio e \(\langle\Omega\rangle\) basso massimizza
\(T_s^{(2),\text{sys}}\). È una posizione structural, non una proprietà
essenziale: qualsiasi sistema può occupare la posizione del Bilanciere
se costruisce la giusta combinazione. Questo è il cuore del programma
politico del libro.



\section{Implicazioni Interpretative}\label{implicazioni-interpretative}

\subsection{Perché la China non vince}\label{perchuxe9-la-cina-non-vince}

Nonostante \(\Phi\) alto e material capacity enorme, la China ha
\(\langle\Omega\rangle\) molto alto (Neijuan structural). Anche con la
divisione \(\langle\Omega\rangle/\Phi\) favorevole, l'involution resta
tossica. Inoltre, l'\(\langle A\rangle\) basso (innovazione vincolata
dalla pianificazione) limita la convessità sistemica. Predizione DEQ$^2$:
\(T_s^{(2),\text{sys}}\) Chinese sopra 1 ma con margine in erosione.

\subsection{Perché gli USA non si
stabilizzano}\label{perchuxe9-gli-usa-non-si-stabilizzano}

L'\(\langle A\rangle\) resta alto ma \(\Phi\) è in declino persistente
(polarizzazione, frammentazione mediatica, secessionismo culturale). Il
prodotto \(\Phi\langle A\rangle\) si comprime; la divisione
\(\langle\Omega\rangle/\Phi\) si amplifica. Predizione DEQ$^2$:
\(T_s^{(2),\text{sys}}\) USA in avvicinamento alla soglia \(\Phi^*\) tra
2027 e 2029.

\subsection{Perché Musk è solitone}\label{perchuxe9-musk-uxe8-solitone}

In quanto attore singolo, Musk non ha un \(\Phi\) definito (\(N=1\)). Il
suo \(A\) individuale è eccezionale, ma non one can aggregare. La sua
\emph{strategia di sopravvivenza} consiste nel restare uno (no
successione, no istituzionalizzazione) --- therefore il suo \(T_s^{(2)}\) è
un caso limite del formalism, non un caso sistemico.

\subsection{Perché il Brazil è candidato a
$\Phi$-engine}\label{perchuxe9-il-brasile-uxe8-candidato-a-ux3c6-engine}

L'Estado Logístico (Cervo) descrive un attore che \emph{esporta
coerenza}: il Brazil non vince per material capacity ma per la sua
capacità di \textbf{proiettare \(\Phi\)} in regioni decoerenti
(Mercosul, BRICS, dialogo nord-sud, mediazione climatica).
L'\(\langle A\rangle\) medio it is sufficient, \(\langle\Omega\rangle\)
contenuto è cruciale. La predizione testabile (Ch.~11): se il Brazil
mantiene \(\Phi > 0{,}7\) tra il 2026 e il 2031,
\(T_s^{(2),\text{sys}}\) Brazilian sale sopra il livello dei tre
concorrenti.



\section{Chapter Summary}\label{sintesi-del-chapter}

\Proved{} \textbf{In una frase:} la stabilità sistemica è la stabilità
local modulata da una \textbf{phase coherence asimmetrica} ---
moltiplicativa per il valore (antifragility), divisoria per lo spreco
(involution) --- che traduce l'asimmetria termodinamica fondamentale
tra ordine e disordine in legge geopolitica.

\[
T_s^{(2),\text{sys}} = \frac{s_{iii} z \delta \cdot (1 + \overbrace{\Phi\langle A\rangle}^{\text{valore amplificato dalla coerenza}})}{\sigma \varepsilon_{\text{dis}} \cdot (1 + \underbrace{\langle\Omega\rangle/\Phi}_{\text{spreco amplificato dalla decoerenza}})}
\]

La forma chiude la catena
\(\Pi(t+1) \to \langle\Pi\rangle_\Phi \to T_s^{(2)} \to T_s^{(2),\text{sys}}\)
aperta nel \S{}3.1bis. Da qui in then, il libro è applicato: il Ch.~5
organizza le quattro variables strutturali
\((\Phi, A, \Omega, \varepsilon_{\text{dis}})\) in quadro coerente per
le applicazioni della Part III.



\section{Closing Epistemic Notes
Chapter}\label{note-epistemiche-di-chiusura-chapter}

{\def\LTcaptype{none} % do not increment counter
{\small\begin{longtable}[]{@{}
  >{\raggedright\arraybackslash}p{(\linewidth - 4\tabcolsep) * \real{0.3333}}
  >{\raggedright\arraybackslash}p{(\linewidth - 4\tabcolsep) * \real{0.3333}}
  >{\raggedright\arraybackslash}p{(\linewidth - 4\tabcolsep) * \real{0.3333}}@{}}
\toprule\noalign{}
\begin{minipage}[b]{\linewidth}\raggedright
\#
\end{minipage} & \begin{minipage}[b]{\linewidth}\raggedright
Claim
\end{minipage} & \begin{minipage}[b]{\linewidth}\raggedright
Status
\end{minipage} \\
\midrule\noalign{}
\endhead
\bottomrule\noalign{}
\endlastfoot
1 & L'hypothesis di mean field di Kuramoto applicata al sistema sociale &
\Hypoth{} working hypothesis di scala \\
2 & La determinazione di \(\Phi\) dalla soglia \(K_c\) & \Proved{}
risultato classico (Kuramoto 1984) \\
3 & La forma \(A_{\text{sys}} = \Phi \langle A\rangle\) & \Hypoth{}
first espansione consistente (alternative non lineari da testare) \\
4 & La forma \(\Omega_{\text{sys}} = \langle\Omega\rangle/\Phi\) &
\Hypoth{} first espansione consistente (alternative non lineari da
testare) \\
5 & L'asimmetria \(\Phi\)-moltiplica-\(A\)
vs.~\(\Phi\)-divide-\(\Omega\) & \Hypoth{} firma termodinamica del
sistema sociale \\
6 & I limiti \(\Phi \to 0, 1\) producono comportamento atteso &
\Proved{} verified analiticamente \\
7 & La soglia \(\Phi^*\) come indicatore precoce & \Hypoth{} hypothesis
operativa testabile \\
8 & La mappatura quattro-attori $\to$ quattro-posizioni & \Hypoth{} quadro
interpretativo, formalizzato in Part III \\
9 & La predizione su \(T_s^{(2),\text{sys}}\) USA in avvicinamento a
\(\Phi^*\) tra 2027-2029 & \Conj{} conjecture datata, falsificabile \\
10 & La predizione su Brazil-$\Phi$-engine condizionata a \(\Phi > 0{,}7\) &
\Conj{} conjecture datata, falsificabile \\
\end{longtable}}
}


